\section{Adversarial Examples and Projected Gradient Methods}
\label{sec:adversarial-examples}
Deep neural networks are known to be vulnerable to small, intentionally crafted perturbations that drastically 
affect model outputs or maximize loss metrics \cite{szegedyIntriguingPropertiesNeural2014,goodfellowExplainingHarnessingAdversarial2015}. 
This vulnerability arises from the fundamental difference between average-case and worst-case generalization: 
While adversarial samples are extremely rare when sampling the input space randomly, they appear in large numbers along certain, 
gradient-guided directions \citeSec{12.4}{aggarwalNeuralNetworksDeep2023}. 
By performing gradient ascent directly on the inputs subject to a bounded $\ell_p$-norm constraint 
$\|\delta\|_p \le \varepsilon$, the approach explicitly targets areas of severe performance degradation. 
In the context of neural control, these adversarial search strategies are adapted from deceptive attacks to serve as 
empirical falsification tools, providing a computationally efficient approach to detect violations of safety or 
stability in a worst-case scenario prior to formal verification \cite{yangLyapunovstableNeuralControl2024}.

In supervised learning, the generation of adversarial examples can be formulated as a constrained optimization problem 
\cite{goodfellowExplainingHarnessingAdversarial2015,madryDeepLearningModels2019}. 
Given a trained model parameterized by $\theta$, a clean input $x \in \mathcal{X}$, 
and its ground-truth target $y$, a white-box adversary searches for an optimal perturbation vector $\delta$ 
that maximizes the model's loss function $\mathcal{L}(x + \delta, y; \theta)$ \cite{madryDeepLearningModels2019}:
\begin{equation}
    \label{eq:adversarial-problem}
    \max_{\delta}
    \quad
    \mathcal{L}(x + \delta, y; \theta)
    \qquad
    \text{s.t.}
    \quad
    \|\delta\|_p \leq \varepsilon,
\end{equation}
where $\varepsilon > 0$ denotes the maximum perturbation magnitude constrained by an $\ell_p$-norm, 
typically chosen as $p \in \{1, 2, \infty\}$ \cite{goodfellowExplainingHarnessingAdversarial2015,madryDeepLearningModels2019}.


\subsection{Projected Gradient Descent}
\label{subsec:pgd}
To solve the inner loss-maximization problem \eqref{eq:adversarial-problem}, the Fast Gradient Sign Method (FGSM) 
\cite{goodfellowExplainingHarnessingAdversarial2015} provides a fast, linear one-step approximation by taking a 
single step along the sign of the loss gradient. However, single-step attacks often fail to find the worst-case loss in complex, 
highly non-convex loss landscapes \cite{madryDeepLearningModels2019}.

Projected Gradient Descent (PGD) extends FGSM into an iterative multi-step gradient ascent scheme \cite{madryDeepLearningModels2019}. 
For an $\ell_\infty$-norm constraint, the admissible perturbation set $\mathcal{C}_{\varepsilon}(x)$ forms a 
hyper-rectangular box centered at $x$:
\begin{equation}
    \mathcal{C}_{\varepsilon}(x)
    =
    \left\{
        z 
        \;\mid\;
        \|z - x\|_p \leq \varepsilon
    \right\}.
\end{equation}
Starting from an initial candidate $x_{\mathrm{adv}} := x$, each step updates the adversarial sample along the gradient direction:
\begin{equation}
    \label{eq:pgd-update}
    x_{\mathrm{adv}}^+
    \leftarrow
    \Pi_{\mathcal{C}_{\varepsilon}(x)}
    \left(
        x_{\mathrm{adv}}
        +
        \alpha \operatorname{sign}
        \left(
            \nabla_{x_{\mathrm{adv}}}
            \mathcal{L}\left(x_{\mathrm{adv}}, y; \theta\right)
        \right)
    \right),
\end{equation}
where $\alpha > 0$ denotes the step size and $\Pi_{\mathcal{C}_{\varepsilon}(x)}(\cdot)$ projects the 
updated sample back onto the constraint set \cite{madryDeepLearningModels2019}. 
Due to the $\ell_\infty$-box structure, this projection simplifies to an element-wise clipping 
operation to the interval $[x - \varepsilon, x + \varepsilon]$, ensuring that $x_{\mathrm{adv}}^+$ 
strictly satisfies the state constraints at every iteration



\subsection{Adaptation to Neural Lyapunov Control}
\label{subsec:pgd-lyapunov-control}